\documentclass[11pt,a4paper]{article}
\usepackage[margin=2.3cm]{geometry}
\usepackage{amsmath,amssymb}
\usepackage{booktabs}
\usepackage[colorlinks=true,linkcolor=black,urlcolor=black]{hyperref}
\pdfstringdefDisableCommands{\def\forall{forall}\def\exists{exists}}

\newcommand{\lab}[1]{a_{#1}}            % label variable
\newcommand{\st}[1]{x_{#1}}             % state variable, start of the step
\newcommand{\stn}[1]{x'_{#1}}           % state variable, end of the step
\newcommand{\SL}{\mathit{SL}}
\newcommand{\R}{R}
\newcommand{\src}{\mathit{src}}
\newcommand{\tgt}{\mathit{tgt}}
\newcommand{\app}{\mathit{app}}
\newcommand{\FA}{\mathit{FA}}
\newcommand{\Bef}{B}

\title{New SAT encodings over factored transition systems}
\author{fd-fts-sat, October 2026}
\date{}

\begin{document}
\maketitle

\section{Notation}

A factored transition system has factors $f \in \mathcal F$ with states $S_f$, labels $L$, and per
factor and label a transition relation $T_f(l) \subseteq S_f \times S_f$. A label is
\emph{irrelevant} in $f$ if $T_f(l) = \{(s,s) \mid s \in S_f\}$; irrelevant labels are skipped in
$f$ everywhere below. We drop the index $f$ where it is clear. For a label $l$ in a factor:
\begin{align*}
  \SL(l) &= \{s \mid (s,s) \in T(l)\} && \text{self-loop states}\\
  \R(l) &= \{(s,s') \in T(l) \mid s \neq s'\} && \text{real transitions}\\
  \app(l) &= \{s \mid \exists s'.\,(s,s') \in T(l)\} && \text{states where $l$ is applicable}\\
  T(l,s) &= \{s' \mid (s,s') \in T(l)\}
\end{align*}
$\prec$ is a total order on $L$ (the label order). Per time step $t$ the encodings use one-hot
state variables $\st{f,s}$ (time $t$) and $\stn{f,s}$ (time $t+1$) with exactly-one constraints
per factor, and label variables $\lab{l}$. Every label fires at most once per step. A step is
satisfiable if the firing labels, executed in the order $\prec$, form a valid sequence from $x$ to
$x'$ under the step semantics of the encoding.

Step semantics, per factor, of the firing labels executed in order $\prec$:
\begin{description}
  \item[$\forall$-step] at most one label moves the factor ($x \to x'$), every other one self-loops at
        both $x$ and $x'$ (in any order).
  \item[$\exists$-step, lax] at most one label moves the factor; a label before the mover self-loops at
        $x$, one after it at $x'$.
  \item[$\exists$-step, strict] lax, and every firing label is applicable at $x$.
  \item[R$^2\exists$] the factor may move several times; each firing label makes one transition from
        the factor's current state (Balyo 2013; the semantics of TAV/TAC and
        \texttt{bdd\_sat(one\_step\_only=false)}).
\end{description}

\section{$\exists$-step on the $\forall$-step encoding (\texttt{label\_sat(exists\_step=true)})}

Base: the $\forall$-step encoding
\begin{center}\texttt{label\_sat(encoding=CHAINS\_PARALLEL, use\_self\_loop\_optimisation=true)}.\end{center}
Its clause families (positive/negative forms and the
row/column/pillar compressions omitted):
\begin{align}
  &\lab{l} \to \textstyle\bigvee_{s \in \SL(l)} \st{s},\quad
   \lab{l} \to \textstyle\bigvee_{s \in \SL(l)} \stn{s} && \R(l) = \emptyset \tag{L1}\\
  &\lab{l} \to \textstyle\bigvee_{s \in \app(l)} \st{s} && \text{rows} \tag{L2}\\
  &\lab{l} \to \textstyle\bigvee_{s' \in \tgt(l)} \stn{s'} && \text{columns} \tag{L3}\\
  &\lab{l} \wedge \st{s} \to \neg\stn{s'} \qquad s' \notin T(l,s) && \text{residual} \tag{L4}\\
  &\text{pillars, self-loop frame} && \text{label-independent} \tag{L5}\\
  &\lab{l} \wedge \stn{s'} \to \neg \lab{l'} \qquad l \neq l' \text{ reach } s',\ s' \notin \SL(l') && \text{chains} \tag{L6}
\end{align}

\paragraph{Chains.} Per factor and target state $t$, the \emph{adders}
$A(t) = (l_1 \prec \dots \prec l_m)$ are the labels with a real transition into $t$. Variables
$I_{t,1},\dots,I_{t,m}$, meaning ``the factor moved into $t$ by $l_k$ or earlier''. With
$\Bef_t(l) = I_{t,k^*}$, $k^* = \max\{k \mid l_k \prec l\}$, and $\Bef_t(l) = \bot$ if there is none:
\begin{align}
  &I_{t,k} \to I_{t,k+1} \tag{E1}\\
  &\lab{l_k} \wedge \st{u} \wedge \stn{t} \to I_{t,k} && (u,t) \in \R(l_k),\ u \notin \SL(l_k) \tag{E2}\\
  &I_{t,k} \to I_{t,k-1} \vee \lab{l_k},\qquad
   I_{t,k} \to I_{t,k-1} \vee \textstyle\bigvee_{(u,t) \in \R(l_k)} \st{u} \tag{E3}\\
  &I_{t,m} \to \stn{t},\qquad I_{t,m} \to \neg\st{t},\qquad \stn{t} \wedge \neg\st{t} \to I_{t,m} \tag{E4}\\
  &\lab{l} \to \neg \Bef_t(l) && t \in \FA(l) \tag{E5}
\end{align}
with $\FA(l) = \{t \mid l \in A(t),\ t \notin \SL(l)\}$ if $\R(l) \neq \emptyset$, and
$\FA(l) = S \setminus \SL(l)$ if $\R(l) = \emptyset$.

\paragraph{Changes to the $\forall$-step clauses (strict).}
(L6) is dropped. (L3) and the $x'$ half of (L1) are dropped for labels with $\SL(l) \neq \emptyset$.
(L4) for $s \neq s'$, $s \in \app(l)$, $s' \notin T(l,s)$ becomes
\[
  \begin{array}{lll}
    s \in \SL(l),\ s' \notin \SL(l): & \lab{l} \wedge \st{s} \wedge \stn{s'} \to \neg \Bef_{s'}(l) & \text{($l$ before the mover)}\\
    s \notin \SL(l),\ s' \in \SL(l): & \lab{l} \wedge \st{s} \wedge \stn{s'} \to \Bef_{s'}(l) & \text{($l$ after the mover)}\\
    s, s' \in \SL(l): & \text{no clause} & \\
    s, s' \notin \SL(l): & \text{unchanged.} &
  \end{array}
\]

\paragraph{Lax.} (L2) and the precondition half of (L1) become, for labels with
$\SL(l) \neq \emptyset$,
\[
  \lab{l} \to \textstyle\bigvee_{s \in \app(l)} \st{s} \vee \bigvee_{s' \in \SL(l)} \Bef_{s'}(l),
  \qquad
  \lab{l} \wedge \st{s} \to \textstyle\bigvee_{s' \in \SL(l) \setminus \{s\}} \Bef_{s'}(l) \quad (s \notin \app(l)).
\]

\paragraph{Semantics.} Exactly $\exists$-step (strict/lax) in the order $\prec$; the mover of a
factor is the adder at which $I_{x',\cdot}$ switches. (E2) omits sources where $l_k$ self-loops, so
such a label may run as a self-loop before the mover; (E3) still lets it be the mover.

\paragraph{Options.}
\begin{description}
  \item[\texttt{exists\_effects\_persist}] keep (L3) for labels with real transitions: a label ends
        the step in one of its targets, as in FDR $\exists$-step.
  \item[\texttt{exists\_pure\_events}] (E2) as $\lab{l_k} \wedge \stn{t} \to I_{t,k}$ for labels with
        $\SL(l_k) = \emptyset$.
  \item[\texttt{exists\_only\_chains}] (diagnostic) (L6) replaced by (E1)--(E5), all other clauses
        unweakened.
\end{description}

\paragraph{Size.} Per step and factor: $\sum_t |A(t)|$ chain variables (one per real
(label, target) pair), $|\R|$ event clauses, $2\sum_t|A(t)|$ upper-bound clauses,
$\sum_l |\FA(l)|$ binary clauses (E5).

\section{Intermediate-state encoding (\texttt{sigma\_sat})}

After Symbolic Pattern Planning (Cardellini, Giunchiglia, Maratea, AIJ 352, 2026). R$^2\exists$
semantics. Per factor and step, the relevant labels in order $l_1 \prec \dots \prec l_m$
(positions). $\sigma^0_s = \st{s}$; $\sigma^i_s$ means ``after the first $i$ positions the factor is
in $s$''. With $C(l) = \{s, s' \mid (s,s') \in \R(l)\}$, position $i$ introduces new variables
$\sigma^i_s$ for $s \in C(l_i)$ only; $\sigma^i_s = \sigma^{i-1}_s$ (same variable) for
$s \notin C(l_i)$. With $l = l_i$, $p = \sigma^{i-1}$, $n = \sigma^i$:
\begin{align}
  &\lab{l} \to \textstyle\bigvee_{s \in \app(l)} p_s \tag{S1}\\
  &\neg\lab{l} \wedge p_s \to n_s,\qquad \neg\lab{l} \wedge n_s \to p_s && s \in C(l) \tag{S2}\\
  &\lab{l} \wedge p_s \to \textstyle\bigvee_{s' \in T(l,s)} n_{s'} && s \in \app(l) \tag{S3}\\
  &\lab{l} \wedge n_{s'} \to \textstyle\bigvee_{s:\,(s,s') \in T(l)} p_s && s' \in C(l) \tag{S4}\\
  &\lab{l} \wedge p_s \to \neg n_{s_1} \vee \neg n_{s_2} && s_1 \neq s_2 \in T(l,s) \tag{S5}\\
  &\sigma^m_s \leftrightarrow \stn{s} \tag{S6}
\end{align}
A label with $\R(l) = \emptyset$ in $f$ gets only (S1). No chain or auxiliary variables.
Variables per step and factor: $\sum_l |C(l)|$; the state of a factor changes only at labels with
real transitions there. Plan extraction: per factor, the state after a fired label is the
$s \in C(l)$ with $\sigma^i_s$ true, else unchanged.

\paragraph{$\exists$-step (\texttt{one\_move=true}, \texttt{strict=true}).} Per factor a monotone
variable $M^i$ for each position with $C(l_i) \neq \emptyset$ (``moved by position $i$''):
\begin{align}
  &M^{i-1} \to M^i,\qquad \neg p_s \wedge n_s \to M^i,\qquad
   M^{i-1} \wedge \neg p_s \to \neg n_s && s \in C(l_i) \tag{M}\\
  &\lab{l} \to \textstyle\bigvee_{s \in \app(l)} \st{s} && \text{strict only} \tag{St}
\end{align}
Semantics: $\exists$-step lax (with (St): strict), in the order $\prec$.

\paragraph{Micro-step groups (\texttt{group\_size=}$k$).} Groups
$G_j = \{l \mid \lfloor \mathrm{pos}_\prec(l)/k \rfloor = j\}$. A group with one relevant label in
$f$ is encoded by (S1)--(S5). A group $G$ with several relevant labels in $f$ is one
$\forall$-step between $p = \sigma^{\mathrm{in}}$ and $n = \sigma^G$, with new variables for
$C(G) = \bigcup_{l \in G} C(l)$ and a mover variable $\mu_l$ per $l \in G$ with $\R(l) \neq \emptyset$:
\begin{align}
  &\lab{l} \to \textstyle\bigvee_{s \in \app(l)} p_s && l \in G \tag{G1}\\
  &\lab{l} \to \textstyle\bigvee_{s \in \SL(l)} n_s && \R(l) = \emptyset \tag{G2}\\
  &\mu_l \to \lab{l} \tag{G3}\\
  &\lab{l} \wedge p_s \to \mu_l && s \in \app(l) \setminus \SL(l) \tag{G4}\\
  &\lab{l} \wedge n_{s'} \to \mu_l && s' \in C(G) \setminus \SL(l) \tag{G5}\\
  &\mu_l \wedge p_s \to \textstyle\bigvee_{s' \in T(l,s)} n_{s'} && s \in \app(l) \tag{G6}\\
  &\mu_l \wedge n_{s'} \to \textstyle\bigvee_{(s,s') \in T(l)} p_s && s' \in C(G) \tag{G7}\\
  &\mu_l \wedge p_s \to \neg n_{s_1} \vee \neg n_{s_2} && s_1 \neq s_2 \in T(l,s) \tag{G8}\\
  &\textstyle\sum_{l \in G} \mu_l \le 1,\qquad
   p_s \wedge \neg n_s \to \bigvee_{l} \mu_l,\qquad \neg p_s \wedge n_s \to \bigvee_l \mu_l
   && s \in C(G) \tag{G9}
\end{align}
Semantics: $\forall$-step inside a group, R$^2\exists$ between groups (sequential in $\prec$).
$k = 1$ is the plain encoding; $k = |L|$ is $\forall$-step.

\paragraph{Per-factor exact blocks (\texttt{factor\_groups=true, factor\_groups\_exact=true}).}
The groups need not be global. In factor $f$, the relevant labels $l_1 \prec \dots \prec l_m$ are
split greedily into blocks of consecutive positions: $l_j$ joins the current block $B$ if
$|B| <$ \texttt{group\_max} and, for every $l \in B$,
\[
  \tgt_f(l) \cap \app_f(l_j) = \emptyset \quad\text{and}\quad \SL_f(l) \cap \src_f(l_j) = \emptyset,
\]
with $\tgt$/$\src$ the targets/sources of real transitions. Under these conditions no two labels
of a block can run one after the other in R$^2\exists$, so a $\forall$-step block allows exactly
the same executions: the horizons equal those of the plain encoding.

\paragraph{Compact block encoding (\texttt{group\_compact=true}).} For a block $G$ with start
$p$, end $n$, new variables $n_s$ for $s \in C(G)$ and $n_s = p_s$ otherwise:
\begin{align}
  &\lab{l} \to \textstyle\bigvee_{s \in \app(l)} p_s && l \in G \tag{C1}\\
  &\lab{l} \wedge p_s \to \textstyle\bigvee_{s' \in T(l,s)} n_{s'} && \R(l) \neq \emptyset \tag{C2}\\
  &\lab{l} \to \textstyle\bigvee_{s \in \SL(l)} n_s && \R(l) = \emptyset \tag{C3}\\
  &p_s \wedge \neg n_s \to \textstyle\bigvee_{l:\, s \in \src(l)} \lab{l},\qquad
   \neg p_s \wedge n_s \to \textstyle\bigvee_{l:\, s \in \tgt(l)} \lab{l} && s \in C(G) \tag{C4}\\
  &\textstyle\sum_{s \in S} n_s \le 1 && \text{ladder} \tag{C5}\\
  &\textstyle\sum_{l \in G,\ \SL(l) = \emptyset,\ \R(l) \neq \emptyset} \lab{l} \le 1 && \text{ladder} \tag{C6}\\
  &\lab{a} \wedge \lab{b} \wedge p_u \wedge n_t \to \bot && (u,t) \in \R(a) \cap \R(b),\ t \notin \SL(a) \cap \SL(b) \tag{C7}
\end{align}
(C5) with (C2)/(C4) makes $n$ one-hot and $(p,n) \in T(l)$ for every firing $l$; (C6) states
directly that labels without self-loops cannot co-fire (C7 is then skipped among them).

\paragraph{Results (43 Autoscale tasks, \texttt{-shr}, algorithm A, goal-chains v2, 300 s).}
\begin{center}
{\small
\begin{tabular}{@{}lrrrrr@{}}
  \toprule
  variant & solved & PAR-2 & time & conflicts & horizon \\
  \midrule
  plain $\sigma$ & 43 & 26.9 & 1.00 & 1.00 & $=$ \\
  fixed groups of 4 & 37 & 117.4 & 0.75 & 0.89 & $+21\%$ \\
  greedy competing groups & 41 & 52.3 & 0.86 & 0.91 & $=$ (one $+1$) \\
  per-factor exact, matrix encoding & 41 & 57.8 & 0.98 & 1.28 & $=$ \\
  per-factor exact, compact, cap 64 & 41 & 49.9 & 0.70 & 0.59 & $=$ \\
  \textbf{per-factor exact, compact + (C6), cap 64} & \textbf{43} & \textbf{25.2} & \textbf{0.72} & \textbf{0.59} & $=$ \\
  \bottomrule
\end{tabular}}
\end{center}
Time and conflicts: geometric means over the tasks all variants solve. Reordering the labels to form
groups (\texttt{group\_order}) changes horizons and was worse in every relaxation tried.

\paragraph{Feeds-DAG reordering (\texttt{commute\_order}).} Label $a$ \emph{feeds} $b$ if in some
factor $f$
\[
  \tgt_f(a) \cap \app_f(b) \neq \emptyset \quad\text{or}\quad \SL_f(a) \cap \src_f(b) \neq \emptyset .
\]
If $a \prec b$ and $a$ does not feed $b$, then $a$ and $b$ fire in one step only as self-loops at a
common state in every shared factor, so they commute. Hence every order $\prec'$ that keeps all
pairs $a \prec b$ with $a$ feeding $b$ admits every one-step label sequence of $\prec$: horizons do
not increase. Modes: 1 ASAP layering of the DAG of forward feeding pairs (level $=1+$ max level of
its feeders), 2 ALAP layering, 3 left/right shift: a label moves before the earliest label it feeds
among those it may pass (no passed label feeds it), symmetrically to the right; every move turns a
backward feeding pair into a forward one, to a fixpoint (deterministic work cap).

Results (43 tasks, as above, deterministic goal-chains order): plain $\sigma$ $\Sigma h = 236$;
modes 1/2 $\Sigma h = 230$ (agricola $12\to11$, scanalyzer $3\to2$, storage $2\to1$,
freecell $5\to4$, driverlog $6\to5$), never higher. ALAP with exact blocks: time 0.75, conflicts 0.51
relative to plain $\sigma$ (exact blocks alone: 0.81, 0.68).

Full Autoscale (386 tasks, all plans valid, no horizon higher): mode 3 with exact blocks solves
331 vs 328 (A, \texttt{-shr}), 335 vs 330 (A, h2+\texttt{mtau}+union), 343 vs 342 (C, \texttt{-shr});
under C, 2 parallel calls keep the coverage with 0.81--0.92 search time, 1 call loses 5--7 tasks.

\paragraph{Lossy blocks (\texttt{group\_relax}).} Ignoring the $\SL \cap \src$ condition
(airport $1 \to 10$) or keeping only goal-chain pairs as block separators (agricola $12 \to 52$)
raises horizons sharply. The feeding pairs that plans use are few (about 4k of 200k adjacent
relations on 43 tasks), but neither goal-chain nor relaxed-plan supporter pairs predict them (about
30\% recall).

\section{Alternating label orders}

\texttt{full\_transitions\_sat(alternate\_label\_orders=true)}: TAV/TAC (R$^2\exists$; transition variables and r2 chains along the label order). Odd time steps
use $\prec$, even ones its reverse. Horizon $2m$ equals Symbolic Pattern Planning with the pattern
$(\prec;\prec^{-1})$ concatenated $m$ times.

\section{Status}

{\small
\begin{tabular}{@{}p{0.3\textwidth}p{0.36\textwidth}p{0.28\textwidth}@{}}
  \toprule
  encoding & branch / commit & validation \\
  \midrule
  \texttt{label\_sat(exists\_step)} & \texttt{exists-speedup} 3a5f6ce85, d150957df, 3fbd3565f &
    horizons $=$ \texttt{exists\_sat}, VAL \\
  \texttt{sigma\_sat} & \texttt{fts-transformations} 7acf888bf (+ uncommitted options) &
    horizons $=$ TAV, \texttt{bdd\_sat} full; VAL \\
  \texttt{sigma\_sat(one\_move, strict)} & \texttt{fts-transformations} 36cbb81a1 & horizons $=$ \texttt{label\_sat}
    $\exists$; VAL \\
  \texttt{sigma\_sat(group\_size)} & \texttt{exists-speedup} 92a8a4221 & VAL; horizon $\ge$ \texttt{group\_size=1} \\
  \texttt{sigma\_sat(factor\_groups, group\_compact)} & \texttt{exists-speedup} f1013138e & VAL; horizons $=$ plain \\
  \texttt{alternate\_label\_orders} & \texttt{fts-transformations} 688d5ba68 & VAL \\
  \bottomrule
\end{tabular}}

\end{document}
